% =====================================================================
% Response to reviewers — [JOURNAL FULL NAME], manuscript [MANUSCRIPT NUMBER]
% Cycle: [second / third / ...] submission, [Month] [Year]
%
% Structure is scaffolding. The replies are the authors' — leave them as
% placeholders unless the user dictates the content.
%
% Rules:
%   - Reviewer comments are quoted VERBATIM. Never paraphrase, soften,
%     or silently correct a reviewer.
%   - Every reply says where the change landed: section, page, equation.
%   - Where a suggestion was not adopted, say so plainly and give the
%     reason. Do not pretend to adopt it.
% =====================================================================

\documentclass[11pt]{article}

\usepackage[T1]{fontenc}
\usepackage[utf8]{inputenc}
\usepackage[sc]{mathpazo}
\usepackage[margin=1in]{geometry}
\usepackage[skip=10pt plus1pt, indent=0pt]{parskip}
\usepackage{amsmath, amssymb}
\usepackage{xcolor}
\usepackage{mdframed}
\usepackage{hyperref}

% Reviewer comment: quoted verbatim, boxed.
\newmdenv[
  linewidth=0.8pt,
  linecolor=gray,
  backgroundcolor=gray!6,
  skipabove=10pt,
  skipbelow=6pt,
  innertopmargin=8pt,
  innerbottommargin=8pt
]{comment}

% Our reply.
\newcommand{\reply}{\medskip\noindent\textbf{Response.}\ }

% Pointer to where the change landed in the manuscript.
\newcommand{\changed}[1]{\medskip\noindent\textit{Change: #1}}

\title{Response to Reviewers\\[4pt]
\large Manuscript [MANUSCRIPT NUMBER]: ``[PAPER TITLE]''\\
\large \emph{[JOURNAL FULL NAME]}}
\author{}
\date{[Month] [Year]}

\begin{document}
\maketitle

\noindent
We thank the reviewers for their careful reading and their comments. We have revised the manuscript
accordingly. Below, each comment is reproduced in full and followed by our response and a pointer to
the corresponding change. Changes in the manuscript are marked in blue.
% ^ delete the last sentence if no marked copy is submitted

\bigskip
\noindent
\textbf{Summary of principal changes}

\begin{itemize}
  \item [CHANGE 1 — what changed, where]
  \item [CHANGE 2]
  \item [CHANGE 3]
\end{itemize}

% =====================================================================
\section*{Reviewer 1}

\begin{comment}
[REVIEWER 1, COMMENT 1 — VERBATIM]
\end{comment}

\reply
[RESPONSE]

\changed{[SECTION / PAGE / EQUATION WHERE THE CHANGE LANDED]}

\begin{comment}
[REVIEWER 1, COMMENT 2 — VERBATIM]
\end{comment}

\reply
[RESPONSE]

\changed{[SECTION / PAGE / EQUATION]}

% =====================================================================
\section*{Reviewer 2}

\begin{comment}
[REVIEWER 2, COMMENT 1 — VERBATIM]
\end{comment}

\reply
[RESPONSE]

\changed{[SECTION / PAGE / EQUATION]}

% ---------------------------------------------------------------------
% Pattern for a suggestion that was NOT adopted. Say so directly.
% ---------------------------------------------------------------------
% \begin{comment}
% [COMMENT — VERBATIM]
% \end{comment}
%
% \reply
% We have not made this change, for the following reason. [REASON]. We hope the reviewer finds this
% acceptable; we have, however, [WHAT WAS DONE INSTEAD, IF ANYTHING] in Section [N].

% =====================================================================
% \section*{Editor}
%
% \begin{comment}
% [EDITOR COMMENT — VERBATIM]
% \end{comment}
%
% \reply
% [RESPONSE]

\end{document}
